\section{Strong typed exact domains on the primitive transversal}
\label{sec:typed-strong}

Throughout this section $U$ is the \emph{one fixed} joint primitive
transversal obtained in Theorem~\ref{thm:joint-transversal}.  Its
finite set of \emph{named} sorts is $I$; in the later EPPA application
these are actual chosen $E$-owners.  Write $L_{ij}$ for the nonempty
finite palette of atomic colours on \emph{distinct} points of sorts
$i,j$.  For $i=j$ the logical diagonal is \emph{not} an element of
$L_{ii}$.  The converse $a\mapsto\bar a$ sends $L_{ij}$ to $L_{ji}$;
all colours in $L_{ii}$ are self-converse.  Every atomic typed triangle
means a triangle on three \emph{pairwise distinct} real points, even
when two or three of the sort labels coincide.  In this section the
symbols $*,\bar A$ always refer to this \emph{strong} typed algebra,
not to raw equality-inclusive relational multiplication in $M$.

The algebraic arguments below use only finite named transitive
sorts, the symmetric complete typed diversity triangle table,
SAP and semantic homomorphic CCO; they do \emph{not} use
primitivity of $U$ or (NR2). Consequently the exact-domain
Helly and finite-meet product results apply to any later
terminal unary-labelled source with those same properties,
\emph{provided} one keeps the strong off-diagonal guard
and does not substitute the equality-inclusive raw product.

For nonempty $A\subseteq L_{ij}$ and $B\subseteq L_{jk}$ define
\[
 A*B=\{c\in L_{ik}:\exists a\in A\,\exists b\in B\;
      (a,b,c)\text{ is a legal typed triangle of types }(i,j,k)\}.
\]
One may also use the empty subset as a temporary Boolean-algebra
zero, but it is never an exact menu.  All products with compatible
nonempty factors are nonempty: strongly amalgamate the two specified
edges over their common endpoint.  By gluing two witnessing triangles
along their common \emph{complete} edge with SAP, with distinct new
endpoints, one obtains
\begin{equation}\label{eq:strong-assoc}
 (A*B)*C=A*(B*C).
\end{equation}
Rotating a legal triangle gives the typed Peirce equivalences
\begin{equation}\label{eq:typed-peirce}
 c\in a*b\ \Longleftrightarrow\ a\in c*\bar b
 \ \Longleftrightarrow\ b\in\bar a*c.
\end{equation}
All of these identities remain true for repeated sort labels; they
never insert an equality edge between distinct formal vertices.

\begin{lemma}[Typed Set-Pasch from original strong amalgamation]
\label{lem:typed-set-pasch}
In any finite complete typed binary SAP source as above,
including the later terminal original-sort source, let
$i,j,k,\ell$ be named sorts, possibly repeated, and let
$A\subseteq L_{ij}$,
$B\subseteq L_{jk}$, $C\subseteq L_{i\ell}$,
$D\subseteq L_{\ell k}$ be nonempty diversity
menus.  Then
\[
 (A*B)\cap(C*D)\ne\varnothing
 \quad\Longrightarrow\quad
 (\bar A*C)\cap(B*\bar D)\ne\varnothing.
\]
No exactness, primitivity, CCO, or identity atom is
required; every displayed product is the genuine STRONG
diversity product between pairwise distinct points.
\end{lemma}
\begin{proof}
Choose a real old $i$--$k$ edge of an atomic colour
$c\in(A*B)\cap(C*D)$.  The two memberships give a
legal old typed triangle $(i,j,k)$ over this edge,
with $i$--$j$ colour in $A$ and $j$--$k$ colour
in $B$, and another legal triangle $(i,\ell,k)$
over the SAME complete $i$--$k$ edge, with
$i$--$\ell$ colour in $C$ and $\ell$--$k$
colour in $D$.  Source SAP amalgamates these
two complete triangles over the whole common edge.
Their two off-base vertices are distinct, even if
$j=\ell$ or other sort labels repeat.  The actual
$j$--$\ell$ diversity atom $e$ on those two vertices
belongs to both $\bar A*C$ (through the old
$i$-vertex) and $B*\bar D$ (through the old
$k$-vertex), proving the required intersection.
\end{proof}

A \emph{positive typed word} is a nonempty sequence of compatible
atomic diversity colours, with arbitrary intermediate sorts.  Its
\emph{strong profile} is its iterated product under $*$; by SAP this
is exactly the set of colours obtainable between the distinct
endpoints of an injectively realised labelled path.  Set
$\mathcal P_{ij}$ to be the collection of these nonempty profiles,
including the full typed top $L_{ij}$ (the empty list of restrictions).
The \emph{strong exact-domain family} is
\begin{equation}\label{eq:typed-H-definition}
 \mathcal H_{ij}=\left\{\bigcap_{r=1}^m P_r\ne\varnothing:
     P_r\in\mathcal P_{ij},\ m\ge0\right\},
 \qquad\bigcap\varnothing=L_{ij}.
\end{equation}
The family of menus is always \emph{tagged} by $(i,j)$; a subset
of $L_{ij}$ is not silently identified with an equal-looking subset
of $L_{k\ell}$.

The description agrees with the guarded \emph{ordinary} path-generated
family used in endpoint semantics.  Here the guard discards equality
and keeps only colours between two distinct formal terminals.  The
next lemma justifies the passage from a represented word, which may
identify nonadjacent path vertices, to a fully injective strong path;
this passage uses SAP of $U$, not SAP of the original raw source $M$.
In particular an empty word or a word consisting only of identity
factors has no guarded diversity survivor and is not inserted as an
artificial strong menu.  Nonempty finite meets commute with this fixed
diversity guard.  Consequently $\mathcal H_{ij}$ is exactly the strong
guard-slice $\mathrm{SEnd}(i,j)$ of the represented endpoint family,
not its artificial unital envelope.

\begin{remark}[The off-diagonal identity firewall]
\label{rem:typed-identity-firewall}
Some ordinary endpoint formalisms adjoin an artificial identity atom
$1_i$ to the loop palette and close under unital products and meets.
The resulting \emph{unital envelope} $\mathrm{OEnd}(i,i)$ need not equal
\eqref{eq:typed-H-definition}.  Its menu $\{1_i\}$ can be artificial
on a pair of distinct formal vertices.  Our $\mathcal H_{ii}$ is the
\emph{strong} family $\mathrm{SEnd}(i,i)$: it consists exclusively of
nonempty \emph{diversity} menus, with no freely adjoined equality.
In particular no equality survivor appears in the arguments below.
For example, a single-sort class with only one diversity colour
$f$ has strong off-diagonal family $\{\{f\}\}$, while its formal
unital envelope also contains the artificial singleton $\{1_i\}$.
\end{remark}

\begin{lemma}[Injectivising typed homomorphisms over fixed roots]
\label{lem:typed-injectivisation}
For this lemma, let $U$ be any homogeneous complete binary relational
structure on named sorts whose age has SAP. No CCO, primitivity or
relative-index assumption is required.
Let $P$ be any finite partial diagram on the named sorts of $U$, with
only diversity atoms on specified edges.  Let $f:P\to U$ be a
colour- and sort-preserving homomorphism, and let $S\subseteq V(P)$
be a set of designated vertices on which $f$ is injective.  For any
finite $Z\subseteq U\setminus f(S)$ there is an \emph{injective}
colour- and sort-preserving homomorphism $f':P\to U$ such that
$f'|_S=f|_S$ and $f'(V(P)\setminus S)\cap Z=\varnothing$.

In particular, for every positive typed word $w$ of diversity atoms,
its represented endpoint relation on \emph{distinct} vertices has
exactly the strong word profile of $w$.  If a represented word includes
identity factors, they may be deleted first; an all-identity word has
an empty diversity guard.  Thus represented multiplication followed
by restriction to a diversity pair palette and strong multiplication
agree on the positive words used in this section.
\end{lemma}
\begin{proof}
First fix any finite complete induced substructure $B\subset U$ and
$v\in U\setminus B$.  Strong amalgamation gives, for every $n$, a
finite amalgam of $n$ copies of $B\cup\{v\}$ over $B$, with the $n$
new points distinct.  Embed it into $U$, then use homogeneity to
identify its copy of $B$ with the prescribed $B$.  Thus
$|\Aut(U/B)\cdot v|$ is infinite.  The argument preserves unary sort
labels, including when the labels of two formal vertices coincide.

Process the finitely many vertices $x\in V(P)\setminus S$ one by one.
Let $B_x$ be the set of the \emph{current} images of vertices joined
to $x$ by a displayed edge.  Since every displayed edge is a
diversity atom, $f(x)\notin B_x$ at that stage.  Its orbit under
$\Aut(U/B_x)$ is infinite, so choose a point of that orbit outside
$Z$ and outside the entire current finite image of $P$ (and, in
particular, outside the fixed image of $S$).  Replace only the image
of the formal vertex $x$ by this point.  All its displayed edge
colours are preserved because their other endpoints are fixed by
the chosen automorphism; all other displayed edges remain unchanged.
The sort of $x$ is also preserved.  The new image of $x$ is distinct
from every other current image, and this property persists as later
vertices are treated.  After all non-designated vertices have been
processed, the resulting map is injective, fixes $S$ and avoids $Z$
outside $S$.

For the word assertion, take any represented realisation of $w$
between two distinct endpoints of the indicated sorts.  Regard the
successive occurrences along the walk as separate formal path
vertices, even if some were sent to the same real point by the
original realisation.  Apply the first assertion with the two
endpoints as $S$.  Their colour is unchanged, and all path vertices
become distinct.  The reverse inclusion follows by forgetting
injectivity.  Logical equality factors may be contracted without
changing the represented endpoints; an all-identity word supplies
only the diagonal, which is excluded by the strong diversity guard.
\end{proof}

\begin{lemma}[SAP-upgraded semantic CCO cycle criterion]
\label{lem:typed-strong-cycle-completion}
Let $U$ be a homogeneous finite-complete named-sorted binary
RELATIONAL structure whose old finite age has SAP and semantic
HOMOMORPHIC CCO. Let $P$ be a finite typed partial diagram in
which every specified pair has ONE legal old DIVERSITY atomic
colour (with its correctly typed converse); all other pairs
are unrecorded. Then $P$ has an INJECTIVE colour-preserving
completion on its original formal vertex set if and only if
EVERY specified simple cycle of $P$ has a colour-preserving
HOMOMORPHISM to $U$.

If $P$ contains a prescribed complete old base $D$, the
injective completion can fix one prescribed actual copy
of $D$. No artificial STRONG identity edges, higher-arity
constraints or partial functions are included.
\end{lemma}
\begin{proof}
An injective completion restricts to a homomorphism on
every specified cycle. Conversely, if every specified
cycle is homomorphically realisable, semantic HOMOMORPHIC
CCO gives ONE colour-preserving homomorphism $f:P\to U$;
it need not be injective. Since all specified pair labels
are genuine irreflexive DIVERSITY atoms, the SAP
injectivisation lemma
\cref{lem:typed-injectivisation} applies to $f$, fixing
any already embedded complete old base $D$, and makes
ALL formal vertices distinct while preserving EVERY
specified typed atomic edge. If a prescribed copy of
$D$ is to be fixed, first use old finite-diagram
homogeneity to align $f(D)$ with it: $f|_D$ is an
induced embedding because $D$ is complete and all its
off-diagonal labels are irreflexive.

The finite induced complete old diagram on the
distinct image points is the desired injective
completion. In particular, by CONTRAPOSITION, absence
of an injective completion forces ONE specified simple
cycle with no colour-preserving homomorphism. This
last inference would be false for merely AP sources:
the repeated-vertex positive word in the following
example has a homomorphic but no injective realisation.
\end{proof}

\begin{example}[The strong guard needs SAP]
\label{ex:typed-walk-SAP-sharp}
The distinction above is real without SAP, even for a symmetric
triangle-defined homogeneous CCO source.  In the four-point
translation structure $\mathbb F_2^2$ of
Example~\ref{ex:exact-ball-sharp}, choose distinct nonzero
$d,e\in\mathbb F_2^2$ and consider the positive diversity word
$w=(a_d,a_d,a_e)$.  Its represented relation between distinct
endpoints is the diversity atom $a_e$: the only walk starting at
$x$ is $x,x+d,x,x+e$.  But this walk repeats its first vertex at
the third position, and no injective realisation of this labelled
path exists.  Thus the guarded represented profile is nonempty
while the corresponding strong profile is empty.  The original
four-point age has AP but not SAP, since two copies of the unique
$a_d$-neighbour of a vertex cannot be strongly amalgamated over
that vertex.  This is \emph{not} a counterexample to the strong
typed calculus, which is invoked only after the full-owner
transversal has been shown to have SAP.
\end{example}

\begin{lemma}[Strong typed theta semantics and $2$-Helly]
\label{lem:typed-helly}
For fixed $(i,j)$, every finite pairwise-intersecting family of
positive typed word profiles has nonempty total intersection.
Consequently every finite pairwise-intersecting family in
$\mathcal H_{ij}$ has nonempty total intersection; a nonempty
intersection of exact menus is exact.  Moreover every
$D\in\mathcal H_{ij}$ is exactly the set of possible atomic colours
between two distinct roots of types $(i,j)$ in a finite strongly
completable typed path diagram.
\end{lemma}
\begin{proof}
For each positive profile $P_r$ choose an actual atomic word defining
it, and connect two \emph{distinct formal roots} of sorts $(i,j)$ by
pairwise internally disjoint paths carrying those words. Before taking the
specified graph, MERGE all duplicate length-one branches: every
length-one profile is one atomic singleton, and pairwise intersection
forces all such singleton atoms to agree. Thus there is at most ONE
literal specified edge between the two roots, rather than parallel
edges or a fictitious two-cycle. If the sole direct atomic
edge is the ONLY root-to-root path in the presentation, the
resulting graph has no specified cycle and its old atom is
already legal; if any longer private path is also present,
the direct edge and that path form a genuine cycle.
Every remaining simple specified cycle is the union
of two different private root-to-root paths (possibly ONE direct
edge together with a longer private path).  The two path profiles meet, so
its edges can be realised by an injective path pair: choose a common
endpoint colour and apply SAP to realise the internal vertices
without identifications.  Thus every simple cycle admits a typed
homomorphism to $U$.  Semantic homomorphic CCO of $U$ completes the
whole theta homomorphically.  The two \emph{formal} roots need not
have distinct images under this particular homomorphism: CCO alone
makes no such promise.

Apply Lemma~\ref{lem:typed-injectivisation} with $S$ consisting
of just \emph{one} root.  Its image stays fixed, but the other root
and all internal vertices are separated, giving a fully injective
completion.  The now-distinct roots have a colour belonging to every
$P_r$.  The use of a one-root pinned set, rather than a two-root
pinned set, is essential when the original CCO map coalesces the
roots.  Only SAP of this \emph{typed} source is used to separate them.

For exact domains, flatten each meet presentation into its list of
positive word profiles.  If two exact menus intersect, every profile
of the first intersects every profile of the second.  Together with
the nonempty individual meets, the combined profile list is
pairwise intersecting, so the preceding argument applies.  Empty
lists present the top and introduce no condition.

To prove the semantic exactness assertion, take a theta presenting
$D=\bigcap_rP_r$, with the bare-root diagram for $D=L_{ij}$.
Every strong realisation gives a root colour in $D$.  Conversely fix
any atom $d\in D$, add the direct root edge labelled $d$, and repeat
the CCO+SAP theta argument: every newly created cycle consists of
the direct edge and one compatible branch.  Thus $d$ is attained.
There is no identity-quotient case because the roots are required
to remain distinct and all $d$ belong to diversity palettes.
\end{proof}

\begin{lemma}[Exact products and finite-meet distributivity]
\label{lem:typed-exact-products}
Every $\mathcal H_{ij}$ contains every singleton atom and its typed
top, is closed under converse and nonempty intersections, and the
family $\mathcal H$ is closed under compatible strong products.
More precisely, for presentations
$A=\bigcap_{r}P_r\in\mathcal H_{ij}$,
$B=\bigcap_{s}Q_s\in\mathcal H_{jk}$, one has
\begin{equation}\label{eq:typed-exact-product}
 A*B=\bigcap_{r,s}(P_r*Q_s),
\end{equation}
with the usual convention that a typed top imposes no restriction.
For a finite nonempty exact meet $B=\bigcap_sB_s$,
\begin{equation}\label{eq:typed-distrib}
 A*B=\bigcap_s(A*B_s),\qquad
 B*Q=\bigcap_s(B_s*Q)
\end{equation}
for every compatible exact $A,Q$.  All displayed products and the
nonempty intersections on the right belong to the indicated exact
families.
\end{lemma}
\begin{proof}
Only the reverse inclusion in \eqref{eq:typed-exact-product} requires
proof.  Let $z$ lie in every $P_r*Q_s$.  First suppose the two
presentations are nonempty.  Choose an atom $b\in B$; by SAP and
Peirce, $z*\bar{\{b\}}$ is nonempty and lies in every
$z*\bar Q_s$.  Thus the profiles $z*\bar Q_s$ meet pairwise,
and the $P_r$ meet pairwise because $A$ is nonempty.  Each cross
intersection $P_r\cap(z*\bar Q_s)$ is nonempty by
$z\in P_r*Q_s$.  The $2$-Helly lemma yields
\[
 a\in A\cap\bigcap_s(z*\bar Q_s).
\]
Now $Q_s\cap(\bar{\{a\}}*\{z\})\ne\varnothing$
for every $s$, while the $Q_s$ share the nonempty menu $B$.
A second application of $2$-Helly gives
$b'\in B\cap(\bar{\{a\}}*\{z\})$, so $z\in a*b'\subseteq A*B$.
The other inclusion follows from monotonicity.  Concatenations
$P_r*Q_s$ are positive word profiles.  When one presentation is empty, use the bare typed top as one
explicit factor in the argument above.  This causes no difficulty:
for all nonempty $P\subseteq L_{jk}$, SAP gives
$L_{ij}*P=L_{ik}$ by amalgamating any chosen $(i,k)$-edge
with a $(j,k)$-edge labelled by an atom of $P$ over its
common $k$-endpoint.  The analogous right-absorption formula
holds by converse.  Thus the empty-presentation convention
agrees with the same product formula.

To prove \eqref{eq:typed-distrib}, expand the exact left multiplier
into its defining positive profiles and use
\eqref{eq:typed-exact-product} twice: both sides are the intersection
of precisely the products of those profiles with the profiles
presenting the $B_s$.  The right-hand identity is symmetric.
No distribution law is asserted for arbitrary, not necessarily
exact, subsets of a palette.
\end{proof}

For $A\in\mathcal H_{ij}$, $B\in\mathcal H_{jk}$ and
$C\in\mathcal H_{ik}$ write
\begin{equation}\label{eq:typed-PC}
 \mathrm{PC}_{ijk}(A,B,C)\quad\Longleftrightarrow\quad
 C\subseteq A*B,\quad
 A\subseteq C*\bar B,\quad
 B\subseteq\bar A*C.
\end{equation}
It requires a legal support triple for \emph{every atom in every
coordinate}.  It is stronger than the existence of one legal triple
inside $A\times B\times C$.  We retain \emph{every} instance,
including repeated sort labels and repeated exact menu variables.
The three conditions transform into one another by Peirce rotation
and converse: in particular
\begin{equation}\label{eq:typed-PC-rotation}
 \mathrm{PC}(A,B,C)\Longleftrightarrow
 \mathrm{PC}(C,\bar B,A)\Longleftrightarrow
 \mathrm{PC}(\bar A,C,B),
\end{equation}
where the endpoint sorts of each rotated triple follow from the
original $(i,j,k)$.

\begin{lemma}[Exact triple occurrence]\label{lem:typed-triple-occurrence}
Three menus $A\in\mathcal H_{ij}$, $B\in\mathcal H_{jk}$,
$C\in\mathcal H_{ik}$ occur simultaneously as the \emph{exact
strong} endpoint domains of three distinct designated vertices in
a finite completable typed partial diagram if and only if
$\mathrm{PC}_{ijk}(A,B,C)$.
\end{lemma}
\begin{proof}
Necessity: if an exact survivor $a\in A$ is realised on the first
pair, its completion has some $b\in B,c\in C$ and those three atoms
form a legal triangle.  Apply the same argument to each coordinate.

Sufficiency: take internally disjoint theta gadgets for the three
menus, one between each pair of three distinct roots of types
$i,j,k$, even if those sort labels coincide. As in
\cref{lem:typed-helly}, merge any duplicate length-one atomic
branches within each root pair; their colours agree since each
menu is nonempty. There are no artificial parallel-edge
two-cycles. Every genuine simple specified cycle lies within
one two-root gadget (and traverses two DIFFERENT branches),
or traverses one branch on each of the three sides.
The first kind is strongly completable by exact-domain
nonemptiness.  For the second, choose any supported legal atomic
triple in $A\times B\times C$; its atoms belong to \emph{every}
branch profile on the corresponding side, so its three labelled
branches can be realised over that terminal triangle by SAP.
Hence all cycles are homomorphically completable.  CCO supplies a
homomorphism, potentially identifying distinct formal roots; apply
Lemma~\ref{lem:typed-injectivisation} with $S=\varnothing$ to obtain
a strong completion of the whole partial diagram.  None of its three
root images is assumed distinct before this final splitting.

Finally, fix any prescribed atom $a$ of one of the three menus.
The appropriate one of the three PC inclusions supplies a supported
legal atomic triangle $(a,b,c)$ with $b\in B,c\in C$.
Add all three corresponding root-edge constraints to the theta
diagram; the preceding cycle argument still applies.  Thus every
atom is feasible on its own designated side.  Since a theta gadget
already restricts that side to its specified intersection, each
exact strong endpoint domain equals the prescribed menu.
\end{proof}
